%=============================================================================
\section{Problem Formulation and Analysis}
\label{sec:method}
%=============================================================================

\subsection{Notation and Problem Formulation}

Let $\Cat = \{c_1,\dots,c_N\}$ be a catalogue and
$\Fset = \{f_1,\dots,f_M\}$ a set of categorical attributes, with $a_f(c)$
denoting the value of attribute $f$ for item $c$. A question is a pair
$q=(f,v)$ with answer function
\begin{equation}
x_q(c) \;=\; \mathbb{1}\!\left[a_f(c)=v\right] \in \{0,1\},
\end{equation}
and $\Qset$ collects all non-degenerate questions. Write $\bx_q \in \{0,1\}^N$
for the column vector of answers and $f(q)$ for the attribute probed by $q$.

The system maintains a belief $\bb_t \in \Delta^{N-1}$, initialised uniform. At
each step it selects $q_t \in \Qset$, observes $y_t \in \{0,1\}$, and updates.
A session runs for at most $T$ questions and terminates with the maximum a
posteriori item. Table~\ref{tab:notation} collects the notation.

\begin{table}[t]
\caption{Notation.}
\label{tab:notation}
\centering
\footnotesize
\begin{tabular}{@{}ll@{}}
\toprule
\textbf{Symbol} & \textbf{Meaning} \\
\midrule
$\Cat$, $N$ & catalogue and its size \\
$\Fset$, $M$ & attribute set and its size \\
$q=(f,v)$, $\Qset$ & question and question space \\
$x_q(c)$ & true answer to $q$ for item $c$ \\
$\bb_t$ & belief over $\Cat$ at step $t$ \\
$c^\ast$ & the user's target item \\
$\varepsilon_f \in [0,\tfrac12)$ & answer error rate for attribute $f$ \\
$\hat\varepsilon_f$ & the value the system assumes \\
$\delta_f \in [0,1)$ & abstention (erasure) rate for attribute $f$ \\
$y_t$ & observed answer at step $t$ \\
$p_q$ & belief-weighted mass answering yes to $q$ \\
$\Hb(\cdot)$ & binary entropy, bits \\
$p \star \varepsilon$ & $p(1-\varepsilon)+(1-p)\varepsilon$ \\
$T$ & query budget \\
\bottomrule
\end{tabular}
\end{table}

\subsection{Observation Model}

We model the user as a noisy channel whose crossover probability depends on the
attribute being probed:
\begin{equation}
y_t \;=\; x_{q_t}(c^\ast) \oplus Z_t, \qquad
Z_t \sim \mathrm{Bernoulli}\!\left(\varepsilon_{f(q_t)}\right),
\label{eq:channel}
\end{equation}
with $Z_t$ independent across turns. The per-attribute index in
\eqref{eq:channel} is the entire departure from the standard formulation, and
everything that follows is a consequence of it.

Three idealisations are worth naming now rather than in the limitations
section. Errors are assumed independent across turns, whereas a user who
misunderstands a category will err consistently within it;
\eqref{eq:channel} admits no abstention, whereas real users decline to answer;
and the channel is symmetric, whereas a user unsure of the answer may be
readier to say yes than no. The first two are not left as caveats:
Section~\ref{sec:erasure} extends the model to abstention and
Section~\ref{sec:results-erasure} measures it, while
Section~\ref{sec:results-correlated} replaces independence with a per-session
absorbing failure and measures that. The third we do not test, and it is listed
again among the threats in Section~\ref{sec:discussion}.

\subsection{Belief Update and a Unification of Elimination Semantics}

Under \eqref{eq:channel} the likelihood of the observed answer is
\begin{equation}
\ell(c \mid q,y) \;=\;
\begin{cases}
1-\hat\varepsilon_{f(q)}, & x_q(c) = y,\\[2pt]
\hat\varepsilon_{f(q)}, & x_q(c) \neq y,
\end{cases}
\end{equation}
and the posterior follows from Bayes' rule,
\begin{equation}
\bb_{t+1}(c) \;=\;
\frac{\bb_t(c)\,\ell(c \mid q_t,y_t)}
     {\sum_{c'\in\Cat} \bb_t(c')\,\ell(c' \mid q_t,y_t)}.
\label{eq:update}
\end{equation}

\begin{proposition}[Elimination as a degenerate limit]
\label{prop:limit}
As $\hat\varepsilon \to 0^{+}$, \eqref{eq:update} converges pointwise to hard
elimination, assigning zero mass to every candidate inconsistent with $y_t$.
At $\hat\varepsilon = 0$ exactly, the normaliser vanishes whenever no candidate
is consistent with the observed history, and the posterior is undefined.
\end{proposition}

\begin{proof}
The likelihood ratio between an inconsistent and a consistent candidate is
$\hat\varepsilon/(1-\hat\varepsilon) \to 0$, giving pointwise convergence. The
normaliser is $\sum_c \bb_t(c)\mathbb{1}[x_q(c)=y]$, which is zero exactly when
the surviving support is empty; since the target is removed from the support by
its first flipped answer, and questions are never repeated, an empty support is
reachable under any positive noise rate.
\end{proof}

Proposition~\ref{prop:limit} has a methodological consequence more important
than its mathematical content. Comparative studies of interactive
identification commonly contrast an ``entropy-based'' system using elimination
against a ``probabilistic'' system using soft updating, and attribute the
resulting difference in noise tolerance to the selection strategy. The two
systems differ in two factors at once, and by Proposition~\ref{prop:limit} the
second factor is not an alternative algorithm but a boundary value of the
first. Any experiment intending to isolate the effect of selection must hold
the update rule fixed. We do so in Section~\ref{sec:baselines}. We also note
that the behaviour of an elimination system on an empty support is an
implementation convention---restart uniformly, retain the prior, or fail---and
that reported accuracies under noise depend on which convention is used, so it
must be stated. Ours is to retain the prior.

\subsection{Reliability-Aware Information Gain}

Let the belief-weighted yes-mass of a question be
\begin{equation}
p_q \;=\; \sum_{c \in \Cat} \bb_t(c)\, x_q(c) \;=\; \bb_t^{\!\top} \bx_q .
\end{equation}
When answers are observed exactly, the expected reduction in belief entropy
from asking $q$ is $\Hb(p_q)$, and maximising it is the classical criterion.
Under \eqref{eq:channel} the system does not observe $x_q(c^\ast)$; it observes
$Y_q$. The relevant quantity is therefore the mutual information between the
target and the observation. Since $Y_q$ is conditionally a BSC output given the
target,
\begin{align}
\RAIG(q)
&= I(C; Y_q) \;=\; H(Y_q) - H(Y_q \mid C) \nonumber\\
&= \Hb\!\left(p_q \star \varepsilon_{f(q)}\right) - \Hb\!\left(\varepsilon_{f(q)}\right).
\label{eq:raig}
\end{align}
The subtracted term is the entropy the channel injects irrespective of the
target and carries no information about it.

We claim no novelty for \eqref{eq:raig} as mathematics: it is the mutual
information of a binary symmetric channel and follows directly from the
observation model~\cite[Ch.~7]{cover2006}. What we claim is that omitting it
is prevailing practice in applied interactive identification, and that the
omission is consequential in a way the next subsections make precise.

\subsection{Analytical Properties}

\begin{proposition}[Consistency and annihilation]
\label{prop:limits}
$\RAIG(q) = \Hb(p_q)$ when $\varepsilon_{f(q)}=0$, and $\RAIG(q)=0$ for every
$p_q$ when $\varepsilon_{f(q)}=\tfrac12$. On $[0,\tfrac12)$, $\RAIG(q)$ is
strictly decreasing in $\varepsilon_{f(q)}$ for $p_q \in (0,1)$.
\end{proposition}

The first limit establishes that classical information gain is the special case
of \eqref{eq:raig} in which the user is assumed infallible, which is the sense
in which the standard criterion is not merely imprecise but makes a specific
and false assumption. The second establishes that a question answered at chance
is worth nothing however well it splits the belief---a property classical
information gain does not possess, since $\Hb(p)$ is blind to $\varepsilon$.

\begin{proposition}[Rank reversal]
\label{prop:reversal}
There exist questions $q_1,q_2$ with $\Hb(p_{q_1}) > \Hb(p_{q_2})$ but
$\RAIG(q_1) < \RAIG(q_2)$.
\end{proposition}

\begin{proof}
Take $p_{q_1}=0.50$, $\varepsilon_{q_1}=0.30$ and $p_{q_2}=0.70$,
$\varepsilon_{q_2}=0.03$, the subjective and objective tier values used
throughout this paper. Then $\Hb(p_{q_1})=1.000 > \Hb(p_{q_2})=0.881$, while
$\RAIG(q_1)=\Hb(0.500)-\Hb(0.30)=1.000-0.881=0.119$ and
$\RAIG(q_2)=\Hb(0.688)-\Hb(0.03)=0.895-0.194=0.701$.
\end{proof}

Proposition~\ref{prop:reversal} is the operative one. Had \eqref{eq:raig}
merely rescaled the classical scores, the argmax would be unchanged and the
correction would be vacuous. It does not: in the instance above the two
criteria disagree by a factor of $5.9$ in the opposite direction, and a system
optimising $\Hb(p)$ will repeatedly spend turns on the question its user can
barely answer. Fig.~\ref{fig:reversal} plots both criteria with this instance
marked.

%---------------------------------------------------------------- Fig. reversal
\begin{figure}[t]
\centering
\begin{tikzpicture}
\begin{axis}[
  width=0.94\columnwidth, height=5.1cm,
  xlabel={split mass $p_q$}, ylabel={score (bits)},
  xmin=0, xmax=1, ymin=0, ymax=1.05,
  xtick={0,0.25,0.5,0.7,1}, xticklabels={0,.25,.5,.7,1},
  ytick={0,0.25,0.5,0.75,1},
  tick label style={font=\scriptsize}, label style={font=\scriptsize},
  legend style={font=\tiny, at={(0.5,1.02)}, anchor=south,
                legend columns=4, draw=none, column sep=3pt},
  grid=major, grid style={black!12, very thin},
  samples=201, domain=0.001:0.999,
]
\addplot[black, thick] {Hb(x)};
\addlegendentry{$\Hb(p)$}
\addplot[blue!70, thick, densely dashed] {RAIGf(x,0.03)};
\addlegendentry{$\varepsilon{=}.03$}
\addplot[orange!85!black, thick, densely dotted] {RAIGf(x,0.15)};
\addlegendentry{$\varepsilon{=}.15$}
\addplot[red!75!black, thick, dashdotted] {RAIGf(x,0.30)};
\addlegendentry{$\varepsilon{=}.30$}
% the rank-reversal instance of Proposition 3
\addplot[only marks, mark=*, mark size=1.6pt, black]
  coordinates {(0.5,1.0) (0.7,0.881)};
\addplot[only marks, mark=square*, mark size=1.7pt, red!75!black]
  coordinates {(0.5,0.119)};
\addplot[only marks, mark=square*, mark size=1.7pt, blue!70]
  coordinates {(0.7,0.701)};
\node[font=\tiny, anchor=south] at (axis cs:0.5,1.0) {$q_1$};
\node[font=\tiny, anchor=south west] at (axis cs:0.7,0.881) {$q_2$};
\draw[->, gray, thin] (axis cs:0.5,0.97) -- (axis cs:0.5,0.16);
\draw[->, gray, thin] (axis cs:0.7,0.855) -- (axis cs:0.7,0.73);
\end{axis}
\end{tikzpicture}
\caption{Classical information gain (solid) against reliability-aware
information gain at three crossover probabilities. Both curves are computed by
the typesetter from \eqref{eq:raig}, not traced. The circles mark the two
questions of Proposition~\ref{prop:reversal} as the classical criterion scores
them, $q_1$ above $q_2$; the squares mark the same two questions once their
attributes' reliabilities are taken into account, and the arrows show how far
each falls. The ordering inverts: $q_1$ is the balanced question on an
attribute answered at $\varepsilon=0.30$, $q_2$ the lopsided one on an
attribute answered at $\varepsilon=0.03$. The $\varepsilon=0.30$ curve never
exceeds $1-\Hb(0.30)=0.119$ bits at any split mass, which is
Corollary~\ref{cor:budget}'s capacity ceiling.}
\label{fig:reversal}
\end{figure}
%----------------------------------------------------------------------

\begin{corollary}[Capacity ceiling and budget bound]
\label{cor:budget}
$\RAIG(q) \le 1 - \Hb(\varepsilon_{f(q)})$, the capacity of the corresponding
channel. Consequently, identifying one of $N$ equiprobable items requires
\begin{equation}
T \;\ge\; \frac{\log_2 N}{\max_{f \in \Fset}\left(1 - \Hb(\varepsilon_f)\right)}
\label{eq:budget}
\end{equation}
questions in expectation.
\end{corollary}

\begin{proof}
The bound on $\RAIG$ is the capacity of BSC$(\varepsilon)$, attained at
$p=\tfrac12$. Identification requires acquiring $\log_2 N$ bits, and no
question delivers more than the largest available capacity; the converse of the
channel coding theorem gives the result~\cite[Ch.~7]{cover2006}.
\end{proof}

Corollary~\ref{cor:budget} is useful in practice rather than deep. It says the
familiar $\log_2 N$ intuition understates the required budget by a factor equal
to the inverse capacity of the most reliable attribute available, and it
supplies a principled way to set $T$ for a given catalogue instead of choosing
it by convention. It is also what fixes the operating point of our experiments:
on the music catalogue it yields $T \ge 20.15$, which is why the primary
comparison is run at $T=20$.

\subsection{Why Discretisation Manufactures the Bias}
\label{sec:discretisation-theory}

The informal claim in Section~I is that quantile binning assigns peak apparent
informativeness to whichever attribute it is applied to, regardless of that
attribute's semantics. That is exactly true, and stating it exactly explains
several otherwise puzzling features of the measurements in
Section~\ref{sec:results-bias}.

\begin{proposition}[Discretisation fixes apparent informativeness]
\label{prop:discretisation}
Let attribute $f$ be obtained by partitioning a continuous variable into $k$
quantile bins of equal population, and let the belief be uniform. Then every
one-versus-rest question on $f$ has split mass exactly $1/k$, so
\begin{equation}
\max_{v}\ \Hb\!\left(p_{(f,v)}\right) \;=\; \Hb(1/k),
\label{eq:binmax}
\end{equation}
which depends on $k$ alone. In particular $k=2$ gives exactly $1$ bit, the
global maximum of $\Hb$, and no attribute of any kind can exceed it. For a
natively categorical attribute with value distribution $\pi$, by contrast,
$\max_v \Hb(\pi_v) = \Hb(1/2)$ only in the coincidental case that some value
carries mass exactly one half, and is strictly smaller otherwise.
\end{proposition}

\begin{proof}
Equal-population quantile binning gives $\pi_v = 1/k$ for every bin $v$, and at
a uniform belief $p_{(f,v)} = \pi_v$. Since $\Hb$ is symmetric about $1/2$ and
increasing on $[0,1/2]$, and $1/k \le 1/2$ for $k \ge 2$, the maximum over
values is attained at any bin and equals $\Hb(1/k)$. The categorical statement
follows because $\Hb$ attains $1$ only at $p=1/2$.
\end{proof}

Two consequences deserve emphasis. First, the apparent informativeness of a
binned attribute is a property of the \emph{binning parameter}, not of the
attribute: danceability and loudness are assigned identical maximum information
gain because both were cut into the same number of bins, not because they are
equally informative about the target. Second, the comparison between a binned
and an unbinned attribute is not a like-for-like comparison of informativeness
at all. A selector maximising $\Hb(p)$ therefore systematically prefers
whichever attributes happened to pass through the discretiser --- and in a
catalogue of perceptual descriptors, those are the subjective ones. This is
what Section~\ref{sec:results-bias} measures, and \eqref{eq:binmax} is why
every four-bin attribute in Table~\ref{tab:attributes} reports the same
$\Hb(1/4)=0.811$ bits.

\begin{corollary}[The cost of $\varepsilon$-blindness]
\label{cor:deficit}
Let $q^\ast = \arg\max_q \Hb(p_q)$ be the question a classical selector asks
and $q^\dagger = \arg\max_q \RAIG(q)$ the reliability-aware choice. The
per-question information deficit incurred by the classical selector is
$D = \RAIG(q^\dagger) - \RAIG(q^\ast) \ge 0$, and
\begin{equation}
D \;\ge\; \Big[\Hb\!\left(p_{q^\dagger} \star \varepsilon_{q^\dagger}\right)
- \Hb(\varepsilon_{q^\dagger})\Big] - \Big[1 - \Hb(\varepsilon_{q^\ast})\Big].
\end{equation}
\end{corollary}

\begin{proof}
Non-negativity is by definition of $q^\dagger$ as the $\RAIG$ maximiser. The
bound substitutes the capacity ceiling of Corollary~\ref{cor:budget} for
$\RAIG(q^\ast)$, which can only decrease the subtrahend.
\end{proof}

Instantiated on the reliability annotation used throughout, a classical
selector that spends a turn on a perfectly balanced subjective question
($\varepsilon = 0.30$) instead of a moderately balanced objective one
($p = 0.7$, $\varepsilon = 0.03$) forgoes at least $0.75 - 0.12 = 0.63$ bits on
that turn. Over a twenty-question budget against an entropy floor of $16.23$
bits, deficits of that size are not a rounding error.

\begin{remark}[Greedy optimality does not transfer automatically]
\label{rem:greedy}
Jedynak et al.~\cite{jedynak2012} establish that greedy one-step entropy
minimisation is optimal over the full horizon for twenty questions with noise,
and that result is often cited as licence to select greedily. It is proved
under a single common crossover probability. We are not aware of an extension
to the heterogeneous per-attribute channel of \eqref{eq:channel}, and we do not
claim one: Corollary~\ref{cor:budget} makes the difficulty visible, since with
heterogeneous capacities a horizon-optimal policy may wish to reserve a
high-capacity attribute for a later turn at which the belief is better
positioned to exploit it, and a myopic rule cannot express that preference. We
select greedily throughout because that is what deployed systems do and because
holding the policy class fixed is what isolates the effect of the acquisition
function, not because optimality has been established in this setting.
\end{remark}

\subsection{Extension: Abstention}
\label{sec:erasure}

Real users decline to answer. Let attribute $f$ additionally carry an
abstention probability $\delta_f$, so that with probability $\delta_f$ the
question returns an erasure and with probability $1-\delta_f$ it returns a
BSC output as before. The channel is then a cascade of a binary symmetric
channel and a binary erasure channel, and because an erasure conveys nothing
about the target while still consuming a turn, the value of the question
becomes
\begin{equation}
\RAIG_\delta(q) \;=\; (1-\delta_{f(q)})
\left[\Hb\!\left(p_q \star \varepsilon_{f(q)}\right) - \Hb(\varepsilon_{f(q)})\right].
\label{eq:erasure}
\end{equation}
The belief update is unchanged on answered turns and is the identity on
erasures. If abstention correlates with subjectivity in the same direction as
error---a user who cannot judge an attribute reliably being also the user most
likely to decline it---then the two penalties compound and the bias should be
\emph{larger} than the pure-BSC analysis suggests. That is a prediction, and
Section~\ref{sec:results-erasure} tests it rather than leaving it as an
expectation.
